\subsection{Safety case: precordial ST from limb leads}
The certificate warns \emph{before} reconstruction that limb-6$\to$precordial is
unsafe for an ST decision: $\kappa{=}\kapLimb$ and the ST segment is only $\dipST$
dipolar, so precordial ST is largely Tier~III. The danger is real and
\emph{bidirectional}. Over $799$ fold-10 records, measuring ST deviation (J$+$60\,ms
vs.\ the PR baseline, $0.1$\,mV threshold) on true vs.\ reconstructed V1--V6, the
generative reconstructor \emph{fabricates} $263$ phantom STEMIs (true signal negative),
whereas the conservative OLS reconstructor \emph{masks} $360$ real ones---opposite,
both-dangerous failures, at an identical $\sim\!0.075$\,mV mean ST error that RMSE
cannot separate. The per-reconstruction flag catches \emph{fabrication} (it detects
added, not missing, energy), abstaining on the highest-$h$ generative ST segments at a
controlled rate; masking is a property of the configuration, which the certificate's
$\kappa$ flags upfront. The lesson is the certificate's, not a single reconstructor's:
this configuration should not drive a precordial ST call. Counts in
\texttt{results/ptbxl\_stemi.json}.

\subsection{Guarantees under device shift}
Calibrating on the Schiller CS device family and testing on the AT family (a genuine
within-PTB-XL device shift), the ST/V4 interval keeps nominal $0.90$ coverage
in-distribution but drops to $0.82$ on the shifted device under plain conformal.
Weighted conformal \cite{tibshirani2019weighted} restores coverage (conservatively, to
$1.00$) and a $300$-record slice recalibration returns it to $0.93$
(\texttt{results/cross\_device.json}). Tier~I exactness and the Tier~III flag's
soundness are unaffected, as promised: only the Tier~II \emph{level} is shift-sensitive,
and it is both measurable and repairable.
